\begin{textsample}{NEG Dim 2 – 2016\_c6 – Score -1.00 – 2016\_c6/enade\_2016\_tec\_gest\_amb\_eac031fe82.txt}  \label{ex:f2_neg_222}
CAPÍTULO 3 – ANÁLISE TÉCNICA DA PROVA Este capítulo tem por objetivo apresentar o desempenho dos estudantes concluintes de Tecnologia em Gestão Ambiental no Enade/2016.
Para isso, foram calculadas as estatísticas básicas da prova em seu todo ( seção 3.1.1 ), bem como as estatísticas relacionadas aos Componentes de Formação Geral ( seção 3.1.2 ) e de Conhecimento Específico da Área ( seção 3.1.3 ).
Dadas as suas características, foram analisadas, em separado, as questões objetivas ( seção 3.2 ) e as questões discursivas ( seção 3.3 ).
Para as questões objetivas, foram disponibilizados os índices de facilidade e de discriminação ponto-bisserial, também em separado, para os Componentes de Formação Geral ( seção 3.2.1 ) e de Conhecimento Específico ( seção 3.2.2 ).
De cada componente, uma das questões foi escolhida para exemplificar a análise gráfica, relacionando as alternativas escolhidas pelos estudantes ( inclusive o gabarito ) com o número de acertos no componente.
O Anexo I apresenta a íntegra da análise gráfica para todas as questões objetivas.
Para cada uma das questões discursivas, os conteúdos dos tipos mais comuns de respostas dos estudantes são apresentados e comparados com o padrão de respostas esperado ( ver Anexo VIII com o padrão de respostas ).
Tomando como base as duas questões discursivas do Componente de Formação Geral, a seção 3.3.1.6 apresenta comentários sobre a correção das respostas com respeito à Língua Portuguesa.
Nas Tabelas 3.1 e 3.3 são apresentados o tamanho da população inscrita, ausente, presente e percentual de faltosos, respectivamente, por Grande Região e por Categoria Administrativa, Organização Acadêmica da IES e Modalidade de Ensino.
As demais tabelas apresentam as seguintes estatísticas das notas15 : média do desempenho na prova, erro padrão da média, desvio padrão, nota mínima, mediana e nota máxima para cada um de seus componentes.
Tais estatísticas contemplam o total de estudantes concluintes da Área de Tecnologia em Gestão Ambiental inscritos e presentes à prova do Enade/2016, tendo em vista agregações, ou por Grandes Regiões e o país como um todo, ou por Categoria Administrativa, Organização Acadêmica da IES e Modalidade de Ensino.
Em relação aos gráficos de distribuição de notas, o intervalo considerado foi de 10 unidades, aberto à esquerda e fechado à direita, com exceção do primeiro intervalo, [ 0 ; 10 ], fechado em ambos os extremos.
Para os gráficos de distribuição das notas das 15 Uma definição dessas estatísticas pode ser encontrada no Glossário. questões discursivas, foram consideradas mais duas categorias : questão em branco16 e nota zero. 3.1 - CONDIÇÃO DE PRESENÇA E ESTATÍSTICAS BÁSICAS DA PROVA Esta seção apresenta a condição de presença, além das estatísticas selecionadas e histogramas da nota geral ( 3.1.1 ) e de cada componente : Formação Geral ( 3.1.2 ) e Conhecimento Específico ( 3.1.3 ).
São também apresentadas estatísticas selecionadas de subpopulações caracterizadas por Grande Região, Categoria Administrativa, Organização Acadêmica e Modalidade de Ensino. 3.1.1 - Condição de Presença e Estatísticas Básicas Gerais A Tabela 3.1 apresenta a Condição de Presença na prova, por Grande Região, dos estudantes concluintes de Tecnologia em Gestão Ambiental.
A população total de Inscritos foi de 6.317.
Destes, 4.833 estiveram Presentes, sendo 23,5 \% o índice de não comparecimento.
A Região de maior ausência foi a Centro-Oeste ( 29,0 \%, ) e a de menor ausência foi a Norte ( 13,4 \% ).
Já a Tabela 3.2 apresenta as Estatísticas Básicas da prova, por Grande Região, dos estudantes concluintes de Tecnologia em Gestão Ambiental.
A Média das notas da prova como um todo ( nas seções seguintes serão analisados os Componentes de Formação Geral e de Conhecimento Específico ) foi 41,3, sendo que os alunos da região Norte obtiveram a Média mais baixa ( 38,1 ), e os da região Nordeste obtiveram a Média mais alta ( 43,0 ).
As demais médias foram : 42,8 na região Sudeste, 40,3 na região Sul e 39,7 na região Centro-Oeste.
O Desvio padrão para o Brasil como um todo foi 12,5, sendo o maior Desvio padrão encontrado na região Sudeste ( 13,0 ) e o menor, na região Centro-Oeste ( 11,4 ), indicando uma dispersão um pouco menor das notas desta última região.
Centro-Oeste ( 11,4 ), indicando uma dispersão um pouco menor das notas desta última região.
As regiões que obtiveram a maior nota Máxima foram a Nordeste e a Sul ( 82,0 ), ao passo que a região que atingiu a menor nota Máxima foi a Centro-Oeste ( 72,0 ).
A Mediana do Brasil como um todo foi 40,9, sendo a maior Mediana obtida na região Sudeste ( 42,6 ), e a menor obtida na Norte ( 37,1 ).
Na região Norte a nota Mínima foi 6,2 e na região Centro-Oeste foi 13,0, nas demais regiões a nota Mínima foi zero.
Considerando -se as notas segundo Grande Região, observa -se que existe diferença estatisticamente significativa ao nível de 95 \% entre a menor Média, obtida na região Norte ( 38,1 ), e as três médias mais altas, obtidas nas regiões Nordeste ( 43,0 ), Sudeste ( 42,8 ) e Sul (40,3)17.
O comportamento das notas dos estudantes de todo o Brasil pode ser observado no Gráfico 3.1 que apresenta um histograma com a distribuição das mesmas.
Essa é uma distribuição unimodal com a moda no intervalo ( 40 ; 50 ], seguido bem de perto pelo intervalo ( 30 ; 40 ]. 17 Todas as comparações deste capítulo utilizam os intervalos de 95 \%.
Os erros-padrão da média que possibilitam os testes estão disponíveis nas tabelas.
A Tabela 3.3 apresenta informações referentes à Condição de Presença desagregadas por Categoria Administrativa, por Organização Acadêmica e por Modalidade de Ensino.
Da população total de Inscritos, 4.667 são de IES Privadas e 1.650 de IES Públicas.
Em relação à Organização Acadêmica, a maior participação foi obtida por estudantes de Universidades ( 3.005 ), seguido de Centros Universitários ( 2.334 ) e Faculdades ( 978 ).
A maioria dos concluintes inscritos na Área de Tecnologia em Educação Ambiental foi da Educação a Distância, 3.561, ao passo que a modalidade presencial contou com 2.756 inscritos.
A Categoria Administrativa de maior ausência foi a Pública ( 23,9 \% ), e entre as Organizações Acadêmicas foi a das Faculdades ( 24,6 \% ), ambas acima da taxa de ausência nacional de 23,5 \%.
Na Modalidade Educação a Distância, a ausência foi de 23,0 \%, e na Educação Presencial, 24,1 \%.
Já a Tabela 3.4 apresenta informações referentes à Média da nota final desagregadas por Categoria Administrativa, por Organização Acadêmica e por Modalidade de Ensino.
A Média das notas da prova como um todo foi 41,3.
Em relação à Categoria Administrativa, os estudantes das IES Públicas obtiveram Média mais alta ( 43,5 ), e os das IES Privadas obtiveram Média mais baixa ( 40,5 ), que a Média nacional.
Observa -se que existe diferença estatisticamente significativa entre as médias das notas das IES Públicas e Privadas.
A diferença entre as médias das regiões Nordeste e Norte ( 4,9 ), a maior e a menor Média, é superior à diferença entre IES Públicas e Privadas ( 3,1 ), caracterizando uma maior diversidade regional do que administrativa.
O Desvio padrão para as IES Públicas ( 12,8 ) foi ligeiramente superior ao do Brasil como um todo ( 12,5 ), indicando uma dispersão um pouco maior das notas nesta Categoria Administrativa.
No tocante à Organização Acadêmica, as Universidades obtiveram Média mais alta do que a nacional ( 42,1 e 41,3, respectivamente ).
A média das Faculdades foi 41,3, igual à Média nacional, e dos Centros Universitários foi menor, 40,2.
Constata -se que existe diferença estatisticamente significativa ao nível de 95 \% entre a maior Média, obtida pelas Universidades, e a Média das notas dos estudantes dos Centros Universitários.
A Média da Modalidade Educação a Distância ( 40,3 ) foi inferior à media da Educação Presencial ( 42,6 ).
Também houve diferença estatisticamente significativa ao nível de 95 \% entre essas duas médias. 3.1.2 - Estatísticas Básicas no Componente de Formação Geral A Tabela 3.5 apresenta as Estatísticas Básicas em relação ao componente da prova que avalia a Formação Geral dos estudantes concluintes.
Os alunos de todo o Brasil obtiveram desempenho médio de 39,9.
Quanto à variabilidade, o Desvio padrão das notas dos estudantes do Brasil, como um todo, foi 15,8.
A maior Média foi obtida na região Nordeste ( 43,6 ), e a menor, na região Sul ( 37,3 ).
As demais médias foram : 39,4 na região Norte, 42,4 na região Sudeste e 38,4 na região Centro-Oeste.
Já o maior Desvio padrão foi obtido na região Sudeste ( 16,0 ), e o menor, na região Norte ( 14,3 ).
Os demais desvios padrões foram : 15,7 nas regiões Nordeste e Sul, e 14,5 da região Centro-Oeste.
A maior nota no Componente de Formação Geral da prova do Enade/2016 foi 95,6, obtida por, pelo menos, um aluno na região Sul.
A menor nota Máxima foi obtida na região Norte ( 74,6 ).
A Mediana do Brasil, como um todo, foi 39,8, sendo a menor Mediana encontrada na região Sul ( 37,5 ), e a maior, encontrada na região Nordeste ( 43,3 ).
A nota Mínima nesta parte foi zero em todas as regiões, sem exceção.
Considerando -se as notas segundo Grande Região, observa -se que existe diferença estatisticamente significativa entre a maior Média das notas do Componente de Formação Geral, obtida na região Nordeste ( 43,6 ), e a Média das regiões Norte, Sul e Centro-Oeste. 3.1.2 - Estatísticas Básicas no Componente de Formação Geral A Tabela 3.5 apresenta as Estatísticas Básicas em relação ao componente da prova que avalia a Formação Geral dos estudantes concluintes.
Os alunos de todo o Brasil obtiveram desempenho médio de 39,9.
Quanto à variabilidade, o Desvio padrão das notas dos estudantes do Brasil, como um todo, foi 15,8.
A maior Média foi obtida na região Nordeste ( 43,6 ), e a menor, na região Sul ( 37,3 ).
As demais médias foram : 39,4 na região Norte, 42,4 na região Sudeste e 38,4 na região Centro-Oeste.
Já o maior Desvio padrão foi obtido na região Sudeste ( 16,0 ), e o menor, na região Norte ( 14,3 ).
Os demais desvios padrões foram : 15,7 nas regiões Nordeste e Sul, e 14,5 da região Centro-Oeste.
A maior nota no Componente de Formação Geral da prova do Enade/2016 foi 95,6, obtida por, pelo menos, um aluno na região Sul.
A menor nota Máxima foi obtida na região Norte ( 74,6 ).
A Mediana do Brasil, como um todo, foi 39,8, sendo a menor Mediana encontrada na região Sul ( 37,5 ), e a maior, encontrada na região Nordeste ( 43,3 ).
A nota Mínima nesta parte foi zero em todas as regiões, sem exceção.
Considerando -se as notas segundo Grande Região, observa -se que existe diferença estatisticamente significativa entre a maior Média das notas do Componente de Formação Geral, obtida na região Nordeste ( 43,6 ), e a Média das regiões Norte, Sul e Centro-Oeste.
O Gráfico 3.2 propicia a avaliação do desempenho dos estudantes do Componente de Formação Geral, a partir do histograma da distribuição das notas correspondentes.
A distribuição é unimodal, com moda em ( 30 ; 40 ], enquanto na prova, como um todo ( Gráfico 3.1 ), a moda foi alcançada no intervalo ( 40 ; 50 ].
Nota -se, ainda que, no Gráfico 3.2, as notas apresentam maior dispersão do que no Gráfico 3.1 ( distribuição das notas da prova ), confirmado pela comparação dos desvios padrões : 12,5 para a nota da prova como um todo e 15,8 para o Componente de Formação Geral.
Na Tabela 3.6 são apresentadas as informações referentes ao desempenho dos concluintes do Componente de Formação Geral, em diferentes agregações : Categoria Administrativa, Organização Acadêmica e Modalidade de Ensino.
Observa -se que existe diferença estatisticamente significativa entre as médias dos tipos de Categoria Administrativa.
A maior Média foi obtida por estudantes de IES Públicas ( 44,1 ), com uma diferença estatisticamente significativa da obtida por estudantes de IES Privadas ( 38,5 ).
Com relação à Organização Acadêmica, nota -se que também há diferença estatisticamente significativa entre as médias dos Centros Universitários e Universidades, respectivamente o menor valor e o maior valor.
No que se refere à Modalidade de Ensino, a Educação Presencial apresentou maior média do que a Educação a Distância ( médias 42,8 e 37,8, respectivamente ).
Tais médias também apresentam diferenças estatisticamente significativas. 3.1.3 - Estatísticas Básicas do Componente de Conhecimento Específico A Tabela 3.7 apresenta as Estatísticas Básicas referentes ao Componente de Conhecimento Específico da área de Tecnologia em Gestão Ambiental.
A Média do desempenho dos alunos do Brasil, como um todo, foi 41,7.
A maior Média foi obtida na região Sudeste ( 42,9 ), e a menor, na região Norte ( 37,7 ).
As demais médias foram : 42,8 na região Nordeste, 41,3 na região Sul e 40,2 na região Centro-Oeste.
Quanto à variabilidade das notas, o Desvio padrão do Brasil, como um todo, foi 13,7, sendo o maior Desvio padrão observado na região Sudeste ( 14,3 ), e o menor, na região Centro-Oeste ( 13,0 ).
Os demais desvios foram : 13,6 na região Norte, 13,1 na região Nordeste e 13,4 na região Sul.
A Mediana das notas dos estudantes de todo o Brasil foi 41,5.
A maior Mediana ocorreu na região Sudeste ( 42,7 ), e a menor, na região Norte ( 37,5 ).
As demais medianas foram : 42,2 na região Nordeste, 41,2 na região Sul e 40,5 na região Centro-Oeste.
A nota Máxima do Brasil, como um todo, foi 95,1, sendo obtida por, pelo menos, um aluno na região Sul.
As demais notas máximas foram : 77,3 na região Norte, 85,4 na região Nordeste, 83,0 na região Sudeste e 77,1 na região Centro-Oeste.
A nota Mínima foi zero nas regiões Nordeste, Sudeste e Sul.
Na região Centro-Oeste foi obtida a maior nota Mínima, 13,6, e na região Norte a nota Mínima foi 3,7.
Observa -se que só não existe diferença estatisticamente significativa da menor Média das notas do Componente de Conhecimento Específico da região Norte ( 37,7 ) e a Média da região Centro-Oeste ( 40,2 ).
Assim como os Gráficos 3.1 e 3.2, o Gráfico 3.3, apresentado a seguir, permite uma avaliação do desempenho de concluintes de Tecnologia em Gestão Ambiental em relação ao Componente de Conhecimento Específico com um histograma da distribuição das notas correspondentes.
Esta também é uma distribuição unimodal, e o grupo modal é o ( 40 ; 50 ], o mesmo grupo modal da prova como um todo e superior ao grupo modal para a Formação Geral.
A Tabela 3.8 apresenta uma comparação dos resultados em relação à Categoria Administrativa, à Organização Acadêmica e a Modalidade de Ensino, agora levando em conta o desempenho de alunos do Componente de Conhecimento Específico da prova.
No que se refere à Organização Acadêmica, a maior Média foi das Universidades ( 42,2 ), vindo a seguir a das Faculdades ( 41,3 ) e, depois, a dos Centros Universitários ( 41,2 ).
O maior Desvio padrão, e acima do valor para o Brasil, como um todo, foi o das Universidades ( 14,1 ).
As Universidades também obtiveram a maior Mediana ( 42,1 ).
Os Centros Universitários obtiveram a maior nota Máxima igual a 95,1.
As Universidades obtiveram nota Máxima igual a 85,4 e as Faculdades obtiveram 79,5.
As demais medianas foram : 41,0 nos Centros Universitários e 41,2 nas Faculdades.
A nota Mínima foi zero para todas as Organizações Acadêmicas.
Observa -se que não existem diferenças estatisticamente significativas ao nível de 95 \% no Componente de Conhecimento Específico entre as notas das Organizações Acadêmicas.
Quanto à Categoria Administrativa, observa -se um comportamento semelhante àquele da parte de Formação Geral e à prova como um todo, ou seja, existe diferença estatisticamente significativa entre as médias das IES Públicas ( 43,3 ) e IES Privadas ( 41,2 ).
Neste caso também, a maior Média foi obtida por alunos de IES Públicas de ensino.
Também semelhante é o comportamento das médias no que toca à Modalidade de Ensino, há diferença estatisticamente significativa entre as médias da Educação Presencial ( 42,5 ) e da Educação a Distância ( 41,1 ). 3.2 - ANÁLISE DAS QUESTÕES OBJETIVAS Esta seção apresenta estatísticas selecionadas e histogramas das Questões Objetivas de Formação Geral ( 3.2.1 ).
São também apresentadas e comparadas as médias das subpopulações caracterizadas por Grande Região. 3.2.1 - Componente de Formação Geral A Tabela 3.9 apresenta as Estatísticas Básicas relativas às oito questões objetivas do componente da prova que abrange a Formação Geral dos estudantes.
A Média do Brasil foi 32,8.
A menor Média foi encontrada na região Sul ( 30,1 ), e a maior, na região Nordeste ( 37,6 ).
As demais médias foram : 30,7 na região Norte, 35,4 na região Sudeste e 30,5 na região Centro-Oeste.
O Desvio padrão do Brasil foi 19,1, sendo o maior Desvio padrão encontrado na região Nordeste ( 20,2 ), e o menor, na região Centro-Oeste ( 16,2 ).
Os demais desvios foram : 17,6 na região Norte, 19,5 na região Sudeste e 18,6 na região Sul.
As medianas do Brasil como um todo e de três das cinco Grandes Regiões foi 25,0 - a Mediana foi maior nas regiões Nordeste e Sudeste, 37,5 em ambas.
A nota Máxima 100,0 foi alcançada por, pelo menos, um estudante das regiões Norte, Sudeste e Sul.
Na região Nordeste a nota Máxima foi 87,5 e na região Centro-Oeste foi 75,0.
As notas mínimas ( 0,0 ) foram iguais para todas as regiões.
A Tabela 3.10 apresenta o Índice de Facilidade e o Índice de Discriminação ( ponto-bisserial ) para cada uma das questões objetivas do Componente de Formação Geral.
Quanto ao Índice de Facilidade, foram usadas as seguintes cores para diferenciar o nível de dificuldade da questão : • Azul para as questões classificadas com índice Muito fácil ( \&gt ; =0,86 ), verde para as questões classificadas com índice Fácil ( 0,61 a 0,85 ), amarelo para as questões classificadas com Médio ( 0,41 a 0,60 ), vermelho para as questões classificadas com Difícil ( 0,16 a 0,40 ) e roxo para as questões classificadas com Muito difícil ( \&lt ; =0,15 ).
Já quanto ao Índice de Discriminação, foram usadas as seguintes cores para qualificar a questão : • As questões classificadas com índice Fraco receberam a cor vermelho ( \&lt ; =0,19 ), as classificadas com Médio receberam a cor amarelo ( 0,20 a 0,29 ), as classificadas com Bom receberam a cor verde ( 0,30 a 0,39 ) e as classificadas com Muito bom ( \&gt ; =0,40 ) receberam a cor azul.
As questões objetivas do Componente de Formação Geral, segundo o Índice de Facilidade, foram assim avaliadas : das oito questões, nenhuma teve o Índice de Facilidade classificado como Muito fácil e uma questão foi tida como Fácil, por ter índice de acertos 0,68.
Uma questão foi considerada com índice de dificuldade Médio, situando -se no intervalo entre 0,41 e 0,60 do Índice de Facilidade, ou seja, houve entre 41,0 \% e 60,0 \% de acertos, enquanto quatro questões foram classificadas na categoria Difícil, situando -se no intervalo entre 0,16 e 0,40.
Por fim, duas questões apresentaram menos de 15 \% de acertos, razão pela qual foram classificadas como Muito difícil.
Como já comentado, para análise das questões objetivas, relativas à Formação Geral, segundo o poder de discriminação, utilizou -se o Índice de Discriminação ( ponto-bisserial ).
Nesta análise, as questões foram assim avaliadas : cinco das oito questões apresentaram índice acima ou igual a 0,40 e, assim, foram classificadas com índice Muito bom para esse grupo de estudantes.
Duas questões tiveram Índice de Discriminação Bom, com valor entre 0,30 e 0,39, e uma questão foi classificada com índice de discriminação Médio ( entre 0,20 e 0,29 ).
Nenhuma questão teve nível Fraco de discriminação para esse grupo de estudantes.
O Índice de Facilidade variou de 0,08 a 0,68, e o de Discriminação, de 0,26 a 0,52.
As cinco questões com Índice de Discriminação Muito bom figuraram entre níveis de dificuldade Fácil, Médio e Difícil : uma classificada na categoria Fácil ( a questões 6 ) do Índice de Facilidade, uma na categoria Médio ( questão 3 ) e três na categoria Difícil ( questões 1, 2 e 7 ).
Em particular, as questões 2 e 3 foram as que apresentaram o maior poder discriminatório, com índice 0,52, sendo a questão 2 considerada Difícil ( índice de facilidade 0,32 ) e a questão 3 considerada Médio ( índice facilidade 0,47 ).
O máximo de acertos foi alcançado pela questão 6 com um Índice de Facilidade de 0,68.
A questão de número 8 apresentou um Índice de Facilidade de 0,17, ou seja, menos do que um quinto dos estudantes conseguiu resolvê -la, dentro do universo de participantes.
Seu Índice de Discriminação foi Bom ( 0,35 ).
Já a questão 7 obteve Índice de Discriminação Muito bom, 0,50, e seu Índice de Facilidade também foi Difícil ( 0,34 ).
O Gráfico 3.4, para exemplificar, analisa o comportamento da questão de número 2 de Formação Geral.
Trata -se de uma questão considerada Difícil, com relação à facilidade e uma das duas que obtiveram o maior Índice de Discriminação dessa parte da prova.
Neste gráfico, cada uma das cinco curvas representa o percentual de respostas em determinada alternativa da questão, em função do número de acertos dos estudantes nessa parte da prova ( Formação Geral/Múltipla Escolha ), antes de possíveis eliminações pelo critério do ponto-bisserial.
A curva em azul corresponde à alternativa A, a alternativa correta para esta questão.
Observa -se que entre os estudantes com menor número de acertos, nessa parte do exame, a situação mais frequente foi a escolha da alternativa E ( em vermelho ), incorreta.
Por exemplo, entre os estudantes que acertaram três questões, 33,2 \% escolheram a alternativa E, 29,0 \% escolheram a alternativa B ( em verde ), 2,2 \% escolheram a alternativa C ( preto ) e 1,6 \%, a D ( roxo ).
Entre os que acertaram três respostas entre as questões de múltipla escolha de formação geral, nenhum deixou a questão em branco e 0,5 \% marcou mais de uma alternativa, invalidando a questão.
À medida em que o número de acertos aumenta, indicando desempenho melhor nessa parte da prova, aumenta concomitantemente a proporção de estudantes que selecionaram a alternativa correta A, atingindo 100 \% para os estudantes com sete acertos.
Essa análise permite verificar como a questão discriminou os grupos de desempenho, justificando o alto índice obtido na questão.
Cumpre notar que não é possível inferir deste gráfico nem o índice de facilidade ( que seria uma média da proporção ponderada pela quantidade de alunos com cada uma das notas ), nem o índice de discriminação ponto-bisserial, por razão equivalente.
No caso extremo no qual a grande concentração dos acertos dos alunos fosse abaixo de 4, o índice de facilidade seria obrigatoriamente abaixo de 40 \% ( neste exemplo ).
Caso a concentração fosse em 6 acertos ou mais, o índice seria obrigatoriamente acima de 60 \%.
Os gráficos relativos às demais questões de Formação Geral constam do Anexo I. 3.2.2 - Componente de Conhecimento Específico A Tabela 3.11 apresenta as Estatísticas Básicas em relação às questões objetivas do Componente de Conhecimento Específico da prova de Tecnologia em Gestão Ambiental por Grande Região.
A Média do Brasil deste componente foi de 42,2.
A menor Média foi observada na região Norte ( 38,0 ), e a maior, na região Sudeste ( 43,3 ).
O Desvio padrão de todo o Brasil foi 14,6, sendo o menor Desvio padrão encontrado na região Centro-Oeste ( 13,8 ), e o maior, na região Sudeste ( 15,0 ).
A Mediana de todo o Brasil foi 43,5, a mesma encontrada em três das regiões : Nordeste, Sudeste e Sul.
Nas regiões Norte e Centro-Oeste a Mediana foi menor, 39,1 em ambas.
A nota Máxima da prova foi 95,7, obtida nas questões objetivas do Componente de Conhecimento Específico, por, pelo menos, um aluno da região Sul.
Nas demais regiões, a nota Máxima da prova foi, em ordem crescente : 78,3 nas regiões Norte e Centro-Oeste, 87,0 na região Sudeste e 91,3 na região Nordeste.
Em todas as regiões a nota Mínima foi zero, exceto na região Centro-Oeste que teve nota Mínima 13,0.
A Tabela 3.12 apresenta os Índices de Facilidade e Discriminação ( ponto-bisserial ) das questões objetivas do Componente de Conhecimento Específico para os estudantes de Tecnologia em Gestão Ambiental.
Para facilitar a diferenciação das questões, usaremos as mesmas cores da Tabela 3.10 para as diferentes classificações dos Índices de Facilidade e de Discriminação.
A partir do Índice de Facilidade obtido, pode -se concluir que um pouco mais de metade das questões objetivas da prova foi considerado, pelo menos, Difícil : das 27 questões válidas, 15 foram classificadas como Difícil ou como Muito difícil.
Três questões foram classificadas como Fácil, e outras nove consideradas como Médio.
Já quanto aos índices de discriminação das questões objetivas do Componente de Conhecimento Específico da prova, tem -se como resultado a seguinte classificação : nove das 27 questões foram consideradas como boas, enquanto uma delas obteve Índice de Discriminação Muito bom.
Assim, para um pouco mais de um terço das questões – 10 em 27 – os Índices de Discriminação foram Bom ou Muito bom.
Dentre as demais, 13 delas foram classificadas como Médio e outras quatro como Fraco, sendo 17, por conseguinte, a quantidade de questões nos dois patamares mais baixos de discriminação.
Constata -se, assim, que a prova – no que se refere ao Componente de Conhecimento Específico – possuía capacidade razoável de discriminar entre aqueles que dominam ou não o conteúdo.
Dentre as questões que alcançaram os maiores índices de discriminação, uma delas, a de número 27, foi classificada com Índice de Discriminação Muito bom, situando -se no intervalo de 0,40 a 0,50 do índice, e foi classificada como Fácil, quanto ao Índice de Facilidade.
A questão de número 28 foi a mais difícil dentre as 27 questões específicas, com baixo Índice de Facilidade, apenas 7 \% de acertos.
Essa questão apresentou poder discriminatório igualmente baixo, -0,02, o que comprova ter sido esta questão a mais difícil para os estudantes.
Destaca -se, também, a questão 14, com Índice de Facilidade 0,11, o que, em termos percentuais, corresponde a 11,0 \% de estudantes que responderam acertadamente.
Já 0,03 foi o seu de Índice de Discriminação.
Tais questões foram, portanto, pelo critério ponto-bisserial, consideradas inadequadas.
Por isso, as questões 14 e 28, além das questões 21 e 32 foram eliminadas do cômputo da nota final.
A título de exemplo das análises do comportamento das questões objetivas, o Gráfico 3.5 analisa a questão 27 do Componente de Conhecimento Específico.
Esta questão foi considerada pelas respostas dos estudantes avaliados na prova como uma questão com índice de facilidade Fácil, 0,67, ou seja, 67 \% dos estudantes assinalaram acertadamente a opção D, correspondente ao gabarito.
Seu Índice de Discriminação foi igual a 0,41, sendo esta questão a que apresentou o maior valor e a única com nível Muito Bom.
Neste gráfico, cada uma das cinco curvas representa o percentual de respostas em determinada alternativa da questão 27, em função do número de acertos dos estudantes nessa parte da prova, antes de possíveis eliminações de questões pelo critério do ponto-bisserial.
A alternativa correta D, representada no gráfico pela curva em roxo, foi escolhida em maiores proporções pelos alunos com desempenho melhor nessa parte da prova.
Já as alternativas incorretas, também denominadas distratores, foram selecionadas, principalmente, por aqueles com notas mais baixas.
Neste caso também a soma não é sempre 100 \% por causa das questões não respondidas ou com mais de uma opção marcada.
Aqueles com nota zero, na sua quase totalidade deixaram esta questão em branco ou marcaram mais de uma alternativa, comportamento considerado inválido.
A proporção de alunos que selecionou a resposta correta D aumenta gradativamente, chegando a atingir 100 \% para 20 acertos, enquanto a proporção dos que escolheram alternativas incorretas decai, a partir de três acertos, como função do número de acertos nessa parte da prova.
Os gráficos relativos às demais questões do Conhecimento Específico constam do Anexo I. 3.3 - ANÁLISE DAS QUESTÕES DISCURSIVAS Esta seção apresenta estatísticas selecionadas e histogramas das Questões Discursivas de Formação Geral ( 3.3.1 ) e Conhecimento Específico ( 3.3.2 ).
São também apresentadas e comparadas as médias de subpopulações caracterizadas por Grande Região. 3.3.1 - Componente de Formação Geral As análises dos resultados de desempenho dos estudantes de Tecnologia em Gestão Ambiental nas duas questões discursivas relativas à Formação Geral encontram -se na Tabela 3.13 e no Gráfico 3.6.
Na Tabela 3.13, observa -se que a nota Média nesse conjunto de questões foi acima da obtida nas objetivas.
Os estudantes, de todo o Brasil, obtiveram, em Formação Geral, Média 32,8 nas questões objetivas e 50,6 nas questões discursivas.
Pode -se notar também que o Desvio padrão nesse conjunto de questões foi maior do obtido nas objetivas, 19,1 nas questões objetivas e 21,9 nas questões discursivas.
A maior Média foi obtida na região Sudeste ( 53,0 ), e a menor, na região Sul ( 48,1 ).
A Mediana de todo o Brasil, neste componente, foi 55,5, a mesma da região Centro-Oeste.
A Mediana foi maior do que a nacional nas regiões Norte ( 57,0 ), Nordeste ( 57,0 ) e Sudeste ( 58,0 ).
Na região Sul foi menor, igual a 53,0.
A nota Máxima ( 94,0 ) foi obtida na região Nordeste.
A menor nota Máxima foi a da região Centro-Oeste ( 85,5 ) e as demais foram : 90,5 na região Norte, 91,5 na região Sudeste e 93,5 na região Sul.
A nota Mínima ( 0,0 ) foi a mesma em todas as regiões do Brasil.
O Gráfico 3.6 representa a distribuição das notas nas questões discursivas do Componente de Formação Geral.
A moda desta distribuição ocorre no intervalo ( 60 ; 70 ].
Destaca -se também o intervalo [ 0 ; 10 ] com distribuição próxima a 8 \% do total de notas, sendo que no intervalo [ 0 ; 10 ] incluem -se além da nota zero, a frequência de alunos que deixaram este tipo de questão em branco.
Na sequência, os resultados verificados para cada uma das questões discursivas de Formação Geral serão apresentados, estabelecendo -se relações com os conteúdos abordados em cada uma delas.
Os comentários da Banca de docentes corretores, a respeito do observado na correção das respostas dos estudantes, suas impressões e conclusões serão apresentados junto à análise de cada questão.
Cumpre esclarecer que, tendo em vista que as questões discursivas de Formação Geral são padronizadas, ou seja, constam de todas as provas, os comentários da Banca são os mesmos para todas as carreiras acadêmicas, sendo direcionados a todos os estudantes que participaram do Enade/2016.
A seguir, serão analisados os desempenhos da Área de Tecnologia em Gestão Ambiental nas duas questões discursivas de Formação Geral do Enade/2016, comparando -se os resultados obtidos com comentários para cada questão. 3.3.1.1 - Análise de Conteúdo da Questão Discursiva 1 do Componente de Formação Geral Os dados de Tecnologia em Gestão Ambiental, obtidos a partir das respostas à questão 1, encontram -se na Tabela 3.14 e no Gráfico 3.7.
Nessa questão - de melhor desempenho dentre as duas de Formação Geral – os alunos, de todo o Brasil, tiveram Média 57,2.
A maior Média para a questão 1 foi obtida nas regiões Norte, Nordeste e Sudeste ( 60,0 ), e a menor, na região Sul ( 54,1 ).
Quanto à variabilidade das notas, o Desvio padrão de todo o Brasil foi 30,8.
O menor Desvio padrão foi obtido na região Norte ( 30,1 ), e o maior Desvio padrão foi obtido na região Centro-Oeste ( 31,9 ).
As medianas das regiões Sul, Centro-Oeste e do Brasil como um todo foram iguais a 60,0.
Nas demais regiões a Mediana foi 70,0.
As notas máximas e mínimas da questão discursiva 1 foram, sem exceção, as mesmas para todas as regiões do Brasil, respectivamente, 100,0 e 0,0.
O Gráfico 3.7 mostra a distribuição das notas na questão discursiva 1 do Componente de Formação Geral.
Observa -se que a maior frequência corresponde aos alunos que obtiveram nota dentro do intervalo ( 70 ; 80 ], correspondendo à moda da distribuição.
Destacam -se, também, os alunos que deixaram a questão em branco com aproximadamente 12 \% do total. 3.3.1.2 - Comentários sobre a correção de Conteúdo das respostas à Questão Discursiva 1 A questão discursiva 1 alinha -se aos conteúdos programáticos previstos nas diretrizes curriculares, quanto às habilidades e às competências para a atualização permanente e aos conhecimentos sobre a realidade brasileira, mundial, tangenciando outras áreas do conhecimento.
O enunciado é claro e indica um tema de extrema relevância, qual seja, as migrações humanas e todo o complexo de problemas que esse movimento causa, quer na origem, quer no seu ponto de destino.
A questão é extremamente atual, maximizada pelos acompanhamentos midiáticos que espetacularizam cenas dramáticas do processo de fuga de massas humanas das guerras, como a morte de imigrantes.
Além do aspecto internacional, deve se ter em conta que o Brasil sofreu vários influxos migratórios, desde o forçado, pela transferência de escravos da África, até o incentivado, dos italianos, alemães e outros europeus para trabalhar na agricultura.
Isso sem considerar a própria colonização portuguesa originária e as migrações posteriores dos mesmos portugueses.
Posteriormente, o Brasil tornou -se um país de emigrantes que foram aportar a outros locais na busca de melhor situação econômica.
Muitos migraram para a América do Norte e para a Europa, principalmente.
Mais recentemente, por força do crescimento econômico pátrio, novas levas de imigrantes desembarcaram no Brasil, muitos deles provenientes do Haiti, embora tenhamos recebidos vizinhos latinos, como bolivianos, paraguaios, chilenos, argentinos, venezuelanos, peruanos, colombianos que também vieram em busca de uma vida melhor.
Como se verifica, a temática indicada no enunciado pertence à realidade conhecida amplamente pelos estudantes, pois, também divulgada por todos os meios de comunicação.
Trata -se de tema geral que, normalmente, não se inclui nas grades curriculares superiores, mas é de conhecimento geral e relacionado aos direitos humanos, este sim, ponto de integração às diretrizes curriculares.
A linguagem utilizada foi clara, indicando a existência dos fatos vinculados aos movimentos migratórios.
Por outro lado, o texto demanda uma reflexão sobre a realidade dos indivíduos que são compelidos ao afastamento dos seus lares por motivos externos, como guerras, conflitos civis, religiosos, perseguições das mais variadas, questões econômicas.
Como se trata de texto contendo uma ampla gama de possibilidades, é natural que ocorram interpretações diversas das esperadas.
Do mesmo modo, a contemporaneidade do tema encaminha o enunciado para ser qualificado como questão de média dificuldade, em comparação com as demais questões apresentadas.
O padrão de resposta indicou possíveis interpretações do texto.
No que concerne ao grau de profundidade, mostrou a necessidade de conhecimento de vários itens pertinentes à migração e que são refletidos nas várias possibilidades de resposta previstas pelo padrão.
Pela análise das provas realizadas, constatou -se que muitos foram os que identificaram os itens indicados no padrão de resposta.
Essa circunstância poderia ser justificada, tendo em vista que o tema vem sendo bastante difundido na mídia social, gerando elementos naturais para a resposta, bem como da história do país, que, intimamente, vincula -se ao processo migratório, quer ativo, remetendo pessoas para outros países, quer passivo, na recepção de pessoas.
Essa característica também permite responder aos termos postos quanto aos Estados que recebem os migrantes.
Assim, a indicação dos aspectos positivos e negativos também é bastante difundida.
As respostas medianas e boas apreenderam o sentido do enunciado, e indicaram os motivos expressos no padrão de respostas, como os ensejadores das imigrações, notadamente, para a Europa, que seria o foco do texto.
Por outro lado, também identificaram os problemas dos imigrantes, no mesmo padrão indicado para as respostas desejadas.
Alguns concluintes procuraram exemplos de comparação na realidade brasileira, como os clássicos imigrantes italianos, alemães e de outras nacionalidades, bem como os atuais haitianos e demais povos da América Latina.
Uma minoria não apreendeu o sentido do enunciado e, somente, apresentou elementos da realidade nacional.
Esses foram avaliados no conjunto de provas fracas.
As respostas demonstraram que o concluinte está atualizado, em regra, com o tema da imigração, indicando os fundamentos dessa movimentação populacional que, nos últimos tempos, tem se agravado pela incidência de conflitos armados em vários países, ensejadores de reflexos na segurança dos habitantes desses locais, bem como o surgimento de focos de doenças e ausência dos serviços públicos básicos, além da fome, decorrente da impossibilidade de produzir alimentos ou mesmo comerciá -los.
Do ponto de vista dos Estados que recebem o fluxo de imigrantes, também restou clara a percepção das dificuldades vinculadas à ausência de estrutura quanto aos serviços públicos ( saúde, educação, segurança, moradia, alimentação ) para receber um número expressivo de imigrantes.
Em época de recessão econômica, a presença de estrangeiros acarreta a competição pela busca de empregos, gera movimentos xenofóbicos, bem como choques culturais e políticos.
Deve -se ter em mira que esses movimentos têm um viés de insegurança, pelos países receptores, diante da crise causada por movimentos considerados terroristas, embora a imensa maioria não tenha qualquer vínculo com esses movimentos.
Os recentes atentados na Europa indicam que os seus autores têm origem em comunidades de imigrantes, embora muitos já tenham nascido nos países da comunidade europeia.
Essas características negativas, no entanto, paradoxalmente, convivem com aspectos positivos, o que encaminha a questão para a necessidade dos jovens 3.3.1.3 - Análise de Conteúdo da Questão Discursiva 2 do Componente de Formação Geral A Tabela 3.15 mostra que o desempenho médio dos estudantes na questão discursiva 2 ( média 41,7 ) foi inferior ao obtido na questão discursiva 1 ( média 57,2 ).
A região Sudeste foi aquela cuja Média foi maior ( 43,7 ), e a de menor Média foi a região Sul ( 39,6 ).
Quanto à variabilidade das notas, o Desvio padrão de todo o Brasil foi 24,6, inferior ao obtido na questão discursiva 1 ( 30,8 ).
O maior desvio nessa questão foi obtido nas regiões Nordeste ( 24,8 ), enquanto o menor foi obtido na região Centro-Oeste ( 23,7 ).
A Mediana de todo o Brasil foi 50,0, a mesma em quase todas as regiões.
Na região Sul, a única exceção, a Mediana foi 40,0.
A região Centro-Oeste foi a única onde a nota Máxima foi menor do que 100,0, igual a 80,0.
As notas mínimas ( 0,0 ) foram as mesmas em todas as regiões, sem exceção.
O Gráfico 3.8 mostra a distribuição das notas na questão discursiva 2 do Componente de Formação Geral.
Observa -se que a maior frequência corresponde aos alunos que obtiveram nota dentro do intervalo ( 50 ; 60 ].
Destacam -se, também, os alunos que deixaram a questão em branco, pouco mais do que 10 \%, caracterizando -se como máximo local.
Nota -se, ainda, que as notas ficaram mais concentradas em imigrantes, que ocuparão postos de trabalho não desejados pelos habitantes locais, bem como, compensarão os baixos índices de crescimento populacional.
Existe outro reflexo positivo que está relacionado ao sistema previdenciário, que terá novos contribuintes nos trabalhadores estrangeiros que se integram ao sistema.
Esses pontos foram abordados e indicam que existe uma consciência dos concluintes dos cursos universitários quanto aos principais itens relacionados ao tema que foi posto para análise. 3.3.1.3 - Análise de Conteúdo da Questão Discursiva 2 do Componente de Formação Geral A Tabela 3.15 mostra que o desempenho médio dos estudantes na questão discursiva 2 ( média 41,7 ) foi inferior ao obtido na questão discursiva 1 ( média 57,2 ).
A região Sudeste foi aquela cuja Média foi maior ( 43,7 ), e a de menor Média foi a região Sul ( 39,6 ).
Quanto à variabilidade das notas, o Desvio padrão de todo o Brasil foi 24,6, inferior ao obtido na questão discursiva 1 ( 30,8 ).
O maior desvio nessa questão foi obtido nas regiões Nordeste ( 24,8 ), enquanto o menor foi obtido na região Centro-Oeste ( 23,7 ).
A Mediana de todo o Brasil foi 50,0, a mesma em quase todas as regiões.
Na região Sul, a única exceção, a Mediana foi 40,0.
A região Centro-Oeste foi a única onde a nota Máxima foi menor do que 100,0, igual a 80,0.
As notas mínimas ( 0,0 ) foram as mesmas em todas as regiões, sem exceção.
O Gráfico 3.8 mostra a distribuição das notas na questão discursiva 2 do Componente de Formação Geral.
Observa -se que a maior frequência corresponde aos alunos que obtiveram nota dentro do intervalo ( 50 ; 60 ].
Destacam -se, também, os alunos que deixaram a questão em branco, pouco mais do que 10 \%, caracterizando -se como máximo local.
Nota -se, ainda, que as notas ficaram mais concentradas em em comparação à questão discursiva de número 1.
Isto pode ser constatado também pela comparação do desvio padrão das notas da questão discursiva 2 ( 24,6 ) e da questão discursiva 1 ( 30,8 ). 3.3.1.4 - Comentários sobre a correção de Conteúdo das respostas à Questão Discursiva 2 A questão discursiva 2 tem seu conteúdo integrado às Diretrizes Curriculares, uma vez que a proteção aos direitos humanos e às relações de gênero constituem matrizes fundamentais para conhecimento dos profissionais de nível superior.
O enunciado é claro e indica um tema de extrema relevância, qual seja, a análise da questão da mulher submetida a injustas agressões e os seus reflexos pessoais, sociais e familiares.
Também enfrenta a busca pela maior proteção no concernente às agressões, com a instituição de um modelo legal mais repressivo e que foi gerado por uma situação especifica de ataque à mulher, com consequências dramáticas para a vítima, Maria da Penha, cuja lei levou seu nome.
Deve ser recordado que os conflitos vinculados às relações familiares sempre foram circunscritos ao seu âmbito e, quando dele extravasavam, por força de agressões contundentes à integridade física de um dos partícipes da relação, em geral a mulher, raramente redundavam em punições relevantes, sendo comum o pagamento em cestas básicas de alimentos, como sanção.
Após a edição da lei Maria da Penha ( Lei no. 11.340, sancionada em 7 de agosto de 2006 ), além das sanções sofrerem modificações, houve a criação de modelos de proteção e atendimento especial para os casos de violência doméstica, procurando proteger as mulheres agredidas antes de danos maiores, muitas vezes irreversíveis, incluindo aí o feminicídio, também incluído no Código Penal, como circunstância qualificadora do crime de homicídio, e incluí -lo no rol dos crimes hediondos.
A lei Maria da Penha busca realizar, como indicado no seu pórtico, a Convenção sobre a Eliminação de Todas as Formas de Discriminação contra as Mulheres e a Convenção Interamericana para Prevenir, Punir e Erradicar a Violência contra a Mulher.
No seu texto, várias são as formas de violência doméstica e familiar contra a mulher, identificadas como sendo, entre outras : “ a violência física, entendida como qualquer conduta que ofenda sua integridade ou saúde corporal ; a violência psicológica, entendida como qualquer conduta que lhe cause dano emocional e diminuição da autoestima ou que lhe prejudique e perturbe o pleno desenvolvimento ou que vise a degradar ou controlar suas ações, comportamentos, crenças e decisões, mediante ameaça, constrangimento, humilhação, manipulação, isolamento, vigilância constante, perseguição contumaz, insulto, chantagem, ridicularização, exploração e limitação do direito de ir e vir ou qualquer outro meio que lhe cause prejuízo à saúde psicológica e à autodeterminação ; a violência sexual, entendida como qualquer conduta que a constranja a presenciar, a manter ou a participar de relação sexual não desejada, mediante intimidação, ameaça, coação ou uso da força ; que a induza a comercializar ou a utilizar, de qualquer modo, a sua sexualidade, que a impeça de usar qualquer método contraceptivo ou que a force ao matrimônio, à gravidez, ao aborto ou à prostituição, mediante coação, chantagem, suborno ou manipulação ; ou que limite ou anule o exercício de seus direitos sexuais e reprodutivos ; a violência patrimonial, entendida como qualquer conduta que configure retenção, subtração, destruição parcial ou total de seus objetos, instrumentos de trabalho, documentos pessoais, bens, valores e direitos ou recursos econômicos, incluindo os destinados a satisfazer suas necessidades ; a violência moral, entendida como qualquer conduta que configure calúnia, difamação ou injúria. ” ( Conforme texto extraído do artigo sétimo da referida lei ).
Também se buscam medidas de integração entre os vários órgãos públicos, abrangendo repressão, proteção, conscientização, bem como utilização de equipes multidisciplinares, tendo em vista os graves danos, de variados matizes, que redundam das práticas agressivas.
Nessa linha, também existe na lei a previsão de locais de atendimento específico ( centros de atendimento integral e multidisciplinar para mulheres e respectivos dependentes em situação de violência doméstica e familiar ; casas-abrigo para mulheres e respectivos dependentes menores em situação de violência doméstica e familiar ; delegacias, núcleos de defensoria pública, serviços de saúde e centros de perícia médico-legal especializados no atendimento à mulher em situação de violência doméstica e familiar ; programas e campanhas de enfrentamento da violência doméstica e familiar ; centros de educação e de reabilitação para os agressores ).
O Poder Judiciário, também, deve especializar o seu atendimento para aplicar as medidas de restrição ao agressor de forma mais eficiente.
A linguagem utilizada na questão foi clara, indicando a grave violação dos direitos das mulheres injustamente agredidas e a luta pela busca de sanções adequadas aos agressores e de proteção pessoal e para os de sua família.
Ademais, o tema não está circunscrito ao Brasil, mas o país possui um número muito expressivo de atos que compõem esse cenário.
Por outro lado, o texto demanda uma reflexão sobre a realidade, e abarca diversos segmentos da sociedade, pois as situações de agressão não se restringem às camadas mais populares.
Como se trata de texto contendo uma ampla gama de possibilidades, é natural que ocorram interpretações diversas das esperadas.
O padrão de resposta identifica a projeção do elaborador na resposta a ser apresentada pelo estudante.
Está em linguagem correta e indica possíveis interpretações do texto.
Suas opções foram bastante amplas, abrangendo a maior parte das respostas encontradas.
No que concerne ao grau de profundidade, o padrão de resposta indicou a necessidade de conhecimento de lei especial que regula a violência doméstica, o que exige ou qualificação da área de Direito ou que os currículos dos cursos dos participantes desse ciclo do Enade tenham oferecido alguma disciplina introdutória da área jurídica.
Ao analisar as provas realizadas, constatamos que, na primeira parte, houve maior facilidade em responder aos itens indicados no padrão de resposta, pois atinentes, não só à formação profissional, mas, também, ao conhecimento da realidade, pois, infelizmente, a violência doméstica está disseminada na sociedade e nos vários pontos do país.
Mas, verificou -se que as respostas vieram, no mais das vezes, desorganizadas, dificultando a correção.
Apesar disso, restou claro que a maioria esmagadora das respostas confluiu para as indicações do padrão apresentado, sendo relevante indicar que o tema remeteu para a luta em relação à igualdade de gênero, temática candente em especial para as concluintes do sexo feminino.
Diversos textos se apresentaram de forma concatenada, fundamentada e, muitas vezes, lamentando que, nos dias atuais, ainda se discuta a necessidade de proteger as mulheres dessa espécie de violência.
Em muitos casos, também, os concluintes fizeram um liame entre o acesso à educação e a melhora das relações entre homens e mulheres, ao permitir que a conclusão de maior grau de educação formal conduzisse à diminuição do índice de violência contra a mulher.
Foram identificadas respostas consoantes com o padrão apresentado, indicando os itens que nele constam em relação aos reflexos pessoais, na família e na sociedade.
As respostas dos estudantes raramente faziam referência à organização protetiva instituída pela Lei Maria da Penha, como esperado pelo padrão de resposta.
A respeito de tal aspecto o padrão indicava como uma das possibilidades de respostas ( item ‘ b ’ ) : a criação de casas de refúgio ou casa-abrigo, para acolher mulheres vítimas de violência doméstica ; reforço às Delegacias de Atendimento à Mulher ; inclusão da vítima em programas assistenciais do governo, programas de proteção à vítima e à testemunha, transferência de local de trabalho, se a vítima for servidora pública.
Como paradigma, indicaram -se as Delegacias da Mulher que são organizadas, com a quase integralidade dos seus membros de policiais do sexo feminino, e especializadas no atendimento a mulheres submetidas a situações de violência.
Essa preocupação gerencial tem por escopo bloquear os preconceitos que acompanham essas dramáticas histórias de violência, o que gera uma incompreensão e uma ineficiência na rede de proteção.
As respostas demonstram que os concluintes têm ciência do problema, da sua gravidade, do vínculo com a questão de gênero e com a submissão da mulher como efeito da sociedade patriarcal, ainda que tratado de forma superficial, com base na mídia e, algumas vezes, com apoio de experiências pessoais.
Identificou -se uma generalizada compreensão de que o problema está sendo enfrentado e que ainda existe um longo caminho no percurso para a sua resolução. 3.3.1.5 - Análise de Língua Portuguesa das Questões Discursivas do Componente de Formação Geral Os dados de Tecnologia em Gestão Ambiental, obtidos a partir das respostas às questões discursivas do Componente de Formação Geral, no que tange à Língua Portuguesa, encontram -se na Tabela 3.16 e no Gráfico 3.9.
Nesse aspecto, os alunos, de todo o Brasil, tiveram Média 55,4.
A maior Média com respeito à Língua Portuguesa foi obtida na região Sudeste ( 57,5 ), e a menor, na região Centro-Oeste ( 53,0 ).
Quanto à variabilidade das notas, o Desvio padrão de todo o Brasil foi 21,1.
O menor Desvio padrão foi obtido na região Norte ( 19,3 ) e o maior Desvio padrão foi obtido na região Sul ( 22,0 ).
A Mediana das notas de Língua Portuguesa foi 60,0 para o Brasil como um todo, a mesma Mediana obtida nas regiões Norte e Sul.
A maior Mediana com respeito à Língua Portuguesa foi obtida nas regiões Nordeste e Sudeste ( 62,5 ), e a menor, na região Centro-Oeste ( 57,5 ).
A nota Máxima para todo o Brasil foi de 92,5, com, pelo menos, um aluno tirando essa nota na região Nordeste.
Além disso, a nota Mínima foi zero em todas as regiões do país.
O Gráfico 3.9 mostra a distribuição das notas de Língua Portuguesa do Componente de Formação Geral.
Observa -se que a maior frequência ( pouco superior a 30 \% ) corresponde aos alunos que obtiveram nota no intervalo ( 60 ; 70 ].
Destaca -se, também, os alunos que deixaram ambas as questões em branco representando em torno de 7,5 \% do total, caracterizando -se como máximo local. 3.3.1.6 - Comentários sobre a correção das respostas de Formação Geral com respeito à Língua Portuguesa Na avaliação do desempenho linguístico dos participantes foram consideradas quatro competências distintas, de modo a permitir um mapeamento detalhado do domínio dos recursos disponíveis na Língua Portuguesa para a comunicação escrita formal.
Nessa situação comunicativa, as exigências para o desenvolvimento desse tipo de texto dizem respeito : à adequação da seleção vocabular, ao desenvolvimento do conteúdo, à estruturação sintática dos períodos, à organização lógica das ideias, à utilização de procedimentos de encadeamento textual e de referenciação, à obediência às exigências morfossintáticas próprias da modalidade escrita da norma-padrão, ao respeito às regras ortográficas e às regras de acentuação gráfica.
A avaliação foi realizada considerando -se os textos das respostas às duas questões de Formação Geral e os resultados obtidos em cada uma delas, como seria de se esperar, foram semelhantes.
As competências avaliadas podem ser sintetizadas nos aspectos descritos a seguir.
Aspectos textuais - estruturação textual condizente com o gênero discursivo e o modo de organização textual expositivo/argumentativo adequado ao gênero – essa competência envolve : a estruturação sintática condizente com o padrão da modalidade escrita formal da Língua Portuguesa de modo a garantir a clareza necessária ; a distribuição do conteúdo do texto em parágrafos, de modo a garantir a sua organização temática ; a utilização de operadores discursivos que contribuam para a progressão temática do texto, estabelecendo relações lógicas entre as ideias apresentadas, tanto do ponto de vista intrafrasal, como do interfrasal ; a utilização de procedimentos de referenciação lexical e pronominal que permitam a retomada de referentes textuais ; o respeito às regras de pontuação como fator de estruturação do período.
Aspectos morfossintáticos - domínio dos diferentes aspectos morfossintáticos próprios da modalidade escrita formal da norma-padrão da Língua Portuguesa – essa competência envolve : a concordância nominal, a concordância verbal, a regência nominal, a regência verbal, a flexão nominal, a flexão verbal, a correlação entre os tempos verbais, a colocação pronominal e a utilização de sinais de pontuação que contribuam para a organização lógica da frase e do texto.
Aspectos vocabulares - seleção vocabular adequada à modalidade escrita formal da Língua Portuguesa, exigida pela situação comunicativa – essa competência envolve a precisão na seleção/utilização do vocabulário relacionado à temática solicitada pela questão ; a ausência de marcas da oralidade, como termos de sentido muito genérico ( “ coisa ”, “ negócio ”, “ você ” ) e termos de registros mais informais ( como gírias, jargões, frases feitas, ditados populares, termos regionais ).
Assim, espera -se que o participante respeite a adequação vocabular não usando gírias ou expressões coloquiais, evite repetição desnecessária de palavras e utilize um vocabulário mais formal, como solicitado por um texto dissertativo.
Aspectos ortográficos - respeito às convenções ortográficas da norma-padrão da Língua Portuguesa – essa competência envolve o domínio das regras de acentuação gráfica e da grafia padrão das palavras ( com ausência de abreviaturas próprias da linguagem da internet ), de acordo com as convenções estabelecidas pela legislação em vigor e consubstanciadas no Vocabulário Ortográfico da Língua Portuguesa, editado pela Academia Brasileira de Letras. ( Vale notar que, nessa edição de 2016, não será aceita a legislação anterior, no caso das regras relativas ao uso do hífen e da acentuação gráfica ).
A escolha dessas competências para subsidiar o processo de avaliação apoia -se na concepção de que, no desempenho dos graduandos, a modalidade escrita tem apresentado uma intensa simplificação, originada no padrão da modalidade oral da Língua Portuguesa.
No caso do texto de base dissertativa, inscrito em um registro formal, a distância entre as duas modalidades é ainda maior, o que provoca situações de hipercorreção ( desvios provocados pela incorporação indevida de uma regra da norma-padrão ) e de truncamentos sintáticos ( estruturas frasais incompreensíveis devido à complexidade sintática própria da modalidade escrita ).
Para efeito de pontuação, o critério aprovado pelo Inep, e já utilizado na avaliação do desempenho linguístico nas edições do Enade desde 2013, reúne as competências acima descritas em três grupos da seguinte forma : aspectos textuais ( 40 \% ), aspectos morfossintáticos e vocabulares ( 40 \% ) e aspectos ortográficos ( 20 \% ).
A seguir, analisam -se qualitativamente os resultados da avaliação de cada um dos aspectos avaliados nos textos de resposta das questões discursivas 1 e 2 de Formação Geral.
Aspectos textuais Esta competência é a que se revelou como a mais problemática entre os participantes, porque são muitos os desvios observados.
Em relação à estrutura formal do texto produzido pelos participantes, observaram -se relações linguísticas quase agramaticais, como as estabelecidas pela sequência de gerúndios sem o apoio de um ponto de partida para a organização das informações gramaticais e semânticas, ou seja, sem oração principal.
Observou -se que uma grande parte dos participantes não organizou as ideias em parágrafos, talvez devido ao pequeno número de linhas disponibilizadas para a resposta da questão ou, quem sabe, pela suposição de que não seria necessária essa divisão por não se tratar de um texto no modelo de uma redação dissertativo-argumentativa, como solicitado nos vestibulares.
Quanto à utilização dos mecanismos de referenciação, deve -se destacar a ocorrência, em uma parte dos textos, de repetições de palavras ou expressões sem a utilização de termos sinônimos ou pronomes, como seria adequado.
Outro aspecto relevante a ser destacado é a quase total ausência de operadores argumentativos, tanto intrafrasais como interfrasais, repercutindo uma tendência atual da mídia escrita.
A respeito do uso dos sinais de pontuação, observaram -se vários desvios nos textos analisados, com a ocorrência, inclusive, de parágrafos sem marca interna de pontuação para separar os períodos.
São os seguintes os tipos de problemas encontrados, que se caracterizam como desvios ao estabelecido pela norma-padrão : a ) uso inadequado de vírgula : • ocorrência entre o sujeito e o predicado, o que é considerado inadequado por interromper a sequência natural da frase ; • ocorrência de apenas uma das vírgulas para intercalar uma palavra, uma expressão ou uma oração encaixada, quando o correto seria a demarcação com uma vírgula antes e uma depois do termo encaixado na frase ; • ocorrência da vírgula no lugar do ponto para separar ideias que constituem períodos distintos ; b ) uso inadequado de ponto e vírgula para separar elementos que deveriam ser separados por vírgula, por estarem dentro do mesmo período ; c ) ausência de ponto final para separar períodos.
Vale observar que não foi penalizada a ausência de vírgula para destacar locuções ou adjuntos adverbiais de pequena extensão deslocados de posição na frase, por ser um uso opcional.
Os problemas relativos à estruturação textual geram, portanto, os seguintes tipos de desvios : • sequência justaposta de ideias sem encaixamentos sintáticos ; • redução de estruturas subordinadas, ao lado do aumento na frequência de estruturas coordenadas e absolutas ; • redução no uso de conectores para expressar relações lógicas essenciais à construção do texto ; • emprego inadequado de operadores que não estabelecem relações lógicas coerentes entre ideias do texto ; • repetição de termos sem a utilização de procedimentos mais sofisticados de substituição ( hiperonímias, hiponímias, nominalizações, expressões metafóricas ) ; • frases fragmentadas que comprometem a estrutura lógico-gramatical ; • frases formadas apenas por oração subordinada, sem oração principal.
Aspectos morfossintáticos e vocabulares Os resultados relativos aos aspectos morfossintáticos ( concordância, regência, colocação de pronomes, flexão nominal e verbal ) evidenciam desvios recorrentes no desempenho dos participantes, relacionados ao seu processo de formação e, principalmente, a hábitos da oralidade, marcados pela informalidade.
Apesar de exemplificarem processos de mudança em curso na Língua Portuguesa, esses desvios foram penalizados por não atenderem às exigências da norma-padrão relativas ao texto escrito em situação formal, como deveria ser o texto dissertativo solicitado no Enade.
O desvio mais frequente, em relação à regência, é a falta do sinal indicativo da crase – isso revela que o usuário não tem consciência de que, sob a forma do termo “ a ”, existe a presença de uma preposição “ a ”, exigida pela regência do termo anterior, combinada a um artigo definido ( por exemplo : “ Devido as grandes guerras que estão acontecendo... ” ) Outros problemas relacionados à regência verbal e à nominal foram identificados em alguns textos dos participantes : ( a ) a ausência de preposição ou seu uso indevido antes de pronome relativo ( por exemplo : “ a região que eles moram não oferece condições de segurança ” ) ; ( b ) o uso indevido do pronome relativo ( por exemplo : “ não pensam na cidade para que eles vão ” ) ; ( c ) utilização de pronome “ onde ” deslocado de sua função locativa ).
Esses processos refletem hábitos da modalidade oral da língua, em situações de registro informal.
Apesar da possibilidade de que essa alteração de regência se generalize no padrão escrito da Língua Portuguesa, como já está ocorrendo até em textos jornalísticos, a ausência da preposição foi penalizada neste processo de avaliação.
Outro desvio observado nos textos diz respeito aos processos de concordância verbal e de concordância nominal.
Quanto à concordância de número, além das ocorrências já conhecidas de falta da marca de plural em verbos ou adjetivos, em função da ocorrência de substantivo plural na função de sujeito, observaram -se dois tipos de desvio já registrados nas pesquisas linguísticas : ( a ) a ausência de marca de plural com sujeito posposto ( por exemplo : “ Foi expulso do país os imigrantes que entraram sem autorização da polícia ” ) ; ( b ) o uso indevido da marca de plural para estabelecer concordância com o plural da locução adjetiva, apesar de o substantivo que funciona como núcleo do sintagma nominal estar no singular ( por exemplo : “ A proteção das mulheres que sofrem violência foram garantidas pela Lei Maria da Penha ” ).
Uma ocorrência generalizada foi a ausência de acento circunflexo na forma plural do presente do indicativo do verbo “ ter ”, que foi considerada como um desvio na concordância verbal e não meramente de acentuação gráfica.
Quanto à concordância de gênero, alguns casos foram observados, normalmente no âmbito de sintagmas nominais longos, em que o adjetivo está afastado do substantivo.
Quanto à questão da colocação pronominal, foram poucos os desvios observados.
Apesar de serem aspectos relacionados à oralidade, observou -se que a maioria dos participantes já incorporou regras como a não introdução da frase por um pronome oblíquo e a próclise na presença de um termo atrator.
Não se adotou, entretanto, o padrão excessivamente formal descrito pelas gramáticas normativas em em relação à posição do pronome oblíquo em locuções verbais, já que esse uso está muito distante da prática cotidiana, até em textos mais formais.
Quanto aos aspectos vocabulares ( precisão semântica do vocabulário utilizado e adequação ao registro formal exigido pelo tipo de texto solicitado no enunciado ), podem ser destacados os seguintes tipos de desvios : expressões da oralidade, seleção vocabular incompatível com o contexto, gerando situações de falta de inteligibilidade ; falta de domínio de vocabulário mais abstrato e de maior complexidade, essencial ao desenvolvimento do texto de base dissertativa.
Não se observou presença de palavras de baixa calão, ao contrário do exame de 2015, em que o tema “ funk ” gerou a utilização desse tipo de vocabulário.
Várias marcas de oralidade foram identificadas, embora não em alta frequência : o uso do pronome relativo “ onde ” como relativo universal, falta de artigo definido antes de substantivo, repetição de palavras por falta de vocabulário, reduções como “ tá ”, “ pra ”, “ pro ”, “ prum ”, expressões informais.
Aspectos ortográficos : O desempenho dos participantes revelou uma diferença muito grande nos dois aspectos de caráter ortográfico analisados neste item : grafia e acentuação.
Observou -se a ocorrência de baixo índice de desvios na grafia das palavras e grande índice de desvios no uso dos sinais de acentuação gráfica.
Em vários casos, ocorre ausência completa de acentuação.
Os resultados revelam que existe uma tendência crescente entre os universitários brasileiros de eliminação da acentuação gráfica, talvez motivada pelos hábitos relacionados às redes sociais ou devido a um conhecimento ainda superficial sobre o Acordo Ortográfico da Língua Portuguesa de 1990.
Os casos mais sistemáticos de eliminação do acento indicador da sílaba tônica são : • palavras proparoxítonas ( por exemplo : “ proximos ”, “ politicos ”, “ publica ”, “ fisica ”, “ psicologica ” ) ; • palavras paroxítonas terminadas em ditongo crescente ( por exemplo : “ necessario ”, “ noticiarios ”, “ individuo ”, “ dependencia ”, “ varios ”, “ propria ”, “ violencia ”, “ ocorrencia ”, “ suicidio ” ) ; • palavras oxítonas ( por exemplo : “ ate ”, “ tambem ”, “ mes ” ) ; • palavras paroxítonas terminadas em ditongo nasal ( por exemplo : “ orgão ”, “ orfão ” ).
Quanto ao domínio das convenções relativas à grafia das palavras, pode -se concluir que os desvios estão majoritariamente relacionados a questões relativas à oralidade : • hipercorreção pela escolha de “ e ” no lugar de “ i ” ou de “ i ” no lugar de “ e ”, por influência de hábitos da oralidade : por exemplo, “ inviado ” ( no lugar de “ enviado ” ), “ entervenção ” ( no lugar de “ intervenção ” ), “ empedimento ” ( no lugar de “ impedimento ) ; • eliminação do “ r ” marcador do infinitivo verbal : por exemplo, “ esta ” ( no lugar de “ estar ” ), “ estuda ” ( no lugar de “ estudar ” ), “ migra ” ( no lugar de “ migrar ” ) ; • utilização de “ s ” no lugar de “ c ” : por exemplo, “ mensionado ” ( no lugar de “ mencionado ” ) ; “ insentivou ” ( no lugar de “ incentivou ” ) ; • utilização de “ s ” no lugar de “ x ” : por exemplo, “ espectativa ” ( no lugar de “ expectativa ”, devido à confusão com o verbo “ esperar ” ) ; e a situação contrária, “ x ” no lugar de “ s ”, “ expantoso ” ( no lugar de “ espantoso ” ) ; • utilização de “ ss ” no lugar de “ ç ” : por exemplo, “ gerassões ” ( no lugar de “ gerações ” ) ; • uso de “ am ” no lugar de “ ão ” na forma de terceira pessoa do singular do presente e do pretérito perfeito do indicativo : por exemplo, “ estam ” ( no lugar de “ estão ” ) ; • utilização de “ l ” no lugar de “ u ” em final de sílaba : por exemplo, “ tralma ” ( no lugar de “ trauma ” ) ; • outros casos de desvios de grafia relacionados à variação diastrática podem ser observados em “ estrupo ” ( no lugar de “ estupro ” ), “ subjulgadas ” ( no lugar de “ subjugadas ” ), “ precoseito ” ( no lugar de “ preconceito ” ), “ soubre ” ( no lugar de “ soube ” ), “ vecendo ” ( no lugar de “ vencendo ” ), “ indesencia ” ( no lugar de “ indecência ” ), “ apolojia ” ( no lugar de “ apologia ” ), “ fulga ” ( no lugar de “ fuga ” ). • inadequação no uso da maiúscula para destacar determinadas palavras-chave do texto, como “ Violência ”, “ Brasileiros ”.
Destaque -se, também, a ocorrência de participantes que grafam os textos inteiramente em caixa alta.
Vale observar que, ao contrário do esperado, não apareceram abreviaturas próprias do “ internetês ”, relacionadas ao uso de redes sociais e de e-mails.
Sua influência, portanto, parece se restringir à ausência dos acentos gráficos.
Em relação aos aspectos textuais, observa -se dificuldade na estruturação textual, o que resulta em truncamentos sintáticos e prejuízo da compreensão do texto, e desrespeito a regras essenciais de pontuação, principalmente pela dificuldade em identificar as relações lógicas que se estabelecem entre as unidades frasais.
Em relação aos aspectos morfossintáticos, domínio das regras de concordância nominal e verbal, mas desrespeito às regras de regência nominal e verbal, principalmente no que diz respeito ao uso do sinal indicativo da crase.
Em relação aos aspectos vocabulares, reduzida utilização de palavras de uso mais formal e abstrato, adequado à construção do texto de base dissertativa.
Em relação às convenções ortográficas, domínio da grafia das palavras, mas falta de atendimento à maioria das regras de acentuação gráfica. 3.3.2 - Componente de Conhecimento Específico Na parte da prova relativa às questões discursivas do Componente de Conhecimento Específico ( Tabela 3.17 ), observa -se que a Média foi bem mais baixa do que para as questões discursivas do Componente de Formação Geral.
Enquanto no Componente de Formação Geral a Média para estudantes de Tecnologia em Gestão Ambiental de todo o Brasil foi 50,6, na parte de Conhecimento Específico a Média foi 38,8.
A maior Média deste componente foi obtida pelos estudantes da região Nordeste ( 42,4 ), e a menor, pelos da região Norte ( 35,8 ).
Quanto à variabilidade das notas, o Desvio padrão de todo o Brasil foi 21,6.
O maior Desvio padrão foi encontrado na região Sudeste ( 22,2 ), e o menor, na região Centro-Oeste ( 20,4 ).
A maior nota Máxima, 100, foi obtida em três regiões : Nordeste, Sudeste e Sul.
Nas demais regiões a nota Máxima foi : 88,3 na região Norte e 95,0 na região Centro-Oeste.
Além disso, a nota Mínima ( 0,0 ) foi obtida por pelo menos um aluno em todas as regiões do Brasil.
A Mediana do Brasil como um todo foi 38,3.
A maior Mediana foi obtida nas regiões Nordeste ( 43,3 ), e a menor, nas regiões Norte e Centro-Oeste ( 36,7 ).
O Gráfico 3.10 representa a distribuição das notas nas questões discursivas no Componente de Conhecimento Específico.
A moda desta distribuição ocorre no intervalo ( 40 ; 50 ].
Destaca -se, também, os alunos que obtiveram nota no intervalo [ 0 ; 10 ] com aproximadamente 12 \% do total.
A análise de cada uma destas questões será feita a seguir. 3.3.2.1 - Análise da Questão Discursiva 3 do Componente de Conhecimento Específico Na questão 3, cujos resultados aferidos encontram -se descritos na Tabela 3.18, a Média dos estudantes de todo o Brasil foi 58,0, o melhor desempenho dentre as três questões discursivas de Conhecimentos Específicos.
A menor Média nessa questão foi obtida pelos alunos das regiões Norte e Centro-Oeste ( 56,1 ), enquanto a maior Média foi obtida na região Nordeste ( 60,4 ).
Quanto à variabilidade das notas, o Desvio padrão de todo o Brasil foi 27,9.
O maior Desvio padrão foi obtido na região Sudeste ( 28,2 ), enquanto o menor foi obtido na região Nordeste ( 26,4 ).
A nota Máxima, 100,0 pontos, foi alcançada por, pelo menos, um aluno de todas as regiões.
A Mediana do Brasil como um todo foi 65,0, a mesma das regiões Nordeste e Sudeste.
Nas demais regiões as medianas foram iguais a 60,0.
A nota Mínima ( 0,0 ) foi obtida por, pelo menos, um estudante em todas as regiões do Brasil.
O Gráfico 3.11 mostra a distribuição das notas na questão discursiva 3, do Componente de Conhecimento Específico da área de Tecnologia em Gestão Ambiental.
Sem levar em consideração que aproximadamente 3 \% dos estudantes deixaram a questão em branco e 8 \% receberam nota zero, essa distribuição tem moda principal no intervalo ( 60 ; 70 ] e um máximo local em ( 80 ; 90 ]. 3.3.2.2 - Comentários sobre as respostas à Questão Discursiva 3 O enunciado da questão foi claro e redigido em linguagem acessível.
Abordava os conceitos de epidemiologia e controle de vetores que fazem parte do campo de conhecimentos de Saúde e Saneamento Ambiental, mais especificamente, na parte de Saúde Pública : controle de vetores, epidemiologia, vigilância e educação em saúde.
Dessa forma, a questão está adequada aos conteúdos previstos nas Diretrizes para Tecnologia em Gestão Ambiental da Portaria Inep nº 302, de 08 de junho de 2016.
A abordagem proposta tinha uma visão sistêmica do ambiente e exigia um raciocínio crítico e criativo na identificação e resolução de problemas, envolvendo competências para planejar, executar e analisar instrumentos de gestão ambiental, bem como planejar e avaliar o uso de tecnologias.
O comando da questão solicitava que o estudante desenvolvesse um texto dissertativo composto de três partes, envolvendo : duas condições que favorecem o desenvolvimento e a proliferação dos mosquitos ; duas medidas de prevenção e controle dos vetores a serem implementadas em residências ; e duas ações governamentais que contribuem para a erradicação das doenças transmitidas pelos mosquitos vetores.
Foi fornecido um texto no enunciado que permitia inferir, por interpretação, uma pequena parte da resposta.
Apesar desse fato, os conhecimentos cobrados eram pouco subjetivos, o que facilitou o estabelecimento de um padrão de resposta que se mostrou adequado para a correção.
Como, anualmente, têm sido realizadas diversas companhas educativas na mídia, relacionadas à prevenção e ao controle da Dengue, que veiculam algumas informações básicas, o nível de exigência da questão pode ser considerado baixo.
Os comandos solicitavam respostas que, mesmo os cidadãos comuns, sem conhecimento dos conteúdos específicos da área, respondessem a parte do que se esperava.
O enunciado estava claro, e o que foi pedido apresentava baixo grau de subjetividade, o que levou a poucas respostas diferentes das estabelecidas pelo padrão de respostas.
Os conhecimentos cobrados foram objetivos, e os alunos, em geral, atenderam ao previsto no padrão.
Alguns parafrasearam o enunciado como numa espécie de introdução da resposta.
O padrão apresentava uma grande abrangência de respostas possíveis, fato que levou à valorização das soluções escolhidas pelos estudantes ao formular suas respostas.
No item ‘ a ’, foram solicitadas duas condições socioambientais que favorecessem o desenvolvimento e a proliferação dos mosquitos vetores.
Nesse item, o padrão garantiu a diversidade encontrada nas respostas dadas pelos estudantes.
Apesar disso, alguns alunos foram pouco específicos ao citarem “ presença de lixo, pneus ou entulho ” ( o padrão esperava como resposta presença de lixo, pneus ou entulho, sem o acondicionamento correto ) ou “ presença de água parada ” ( o padrão esperava como resposta água parada em regiões impermeáveis sem o devido acondicionamento, tratamento ou limpeza ).
As respostas com imprecisões destes tipos foram parcialmente pontuadas.
No item ‘ b ’, foram solicitadas duas medidas de prevenção e controle dos vetores a serem implantadas em residências.
A maioria dos alunos apresentou um dos exemplos dentre aqueles presentes no padrão de resposta.
Alguns alunos citaram medidas de controle individual, como uso de repelentes corporais, e outros citaram ações governamentais que não são de uso residencial.
No item ‘ c ’, foram solicitadas duas ações governamentais que contribuem para erradicação de doenças transmitidas pelos mosquitos vetores.
A maioria dos alunos apresentou um dos exemplos do padrão de resposta.
No entanto, houve alunos que, no lugar de propor ações, redigiram suas respostas apontando deficiências nas ações governamentais como : “ falta de saneamento básico ” ou “ falta de informação.
Quanto à linguagem, de forma geral, os alunos apresentaram muitos erros de português e nem sempre utilizaram termos específicos da Área.
Observou -se uma maior aderência ao jargão específico da área nos alunos com melhor desempenho na questão.
De forma geral, pode -se avaliar o desempenho dos alunos como de mediano a bom, fato explicado pelo nível de facilidade da questão, abrangendo tema amplamente divulgado por campanhas educativas da mídia, relativas à prevenção e ao controle da Dengue.
O número de respostas em branco, nulas ou desconsideradas, assim como de notas zero foi reduzido ( menor que 1 \% ), valor bem mais baixo do que em edições anteriores do Enade.
Este fato também sugere um alto grau de comprometimento dos alunos com o exame e com sua universidade.
Nas respostas consideradas fracas, em geral, o estudante confundiu conceitos ou os descreveu de forma trocada, por exemplo, citando ações governamentais na parte das medidas de prevenção e controle em residências.
Houve, também, respostas nas quais o texto era um amontoado de palavras desconectadas ou que fugiram, completamente, do que se pretendia avaliar.
Houve respostas genéricas, como as que citavam falta de infraestrutura como condição socioambiental.
Nas respostas consideradas medianas, em geral, o estudante demonstrou conhecer, parcialmente, o solicitado no enunciado.
Em muitos casos, o estudante respondeu, apenas, a dois itens da questão deixando, muitas vezes, um terceiro sem ser totalmente respondido.
Nas respostas consideradas boas, em geral, o estudante demonstrou conhecer o assunto, articulou os itens, lidou adequadamente com o jargão da área, desenvolveu texto organizado, coerente e conceitualmente correto, e os possíveis deslizes não comprometeram a solução apresentada.
Não houve muita diversidade nas soluções apresentadas, e o padrão foi capaz de capturar as variações corretas de resposta.
Na resposta ao item ‘ a ’, a maioria dos alunos apresentou um dos exemplos do padrão de resposta, de condições que favorecem o desenvolvimento e a proliferação dos mosquitos vetores.
Foram muito comuns as respostas “ deficiência de saneamento básico ”, “ presença de lixo a céu aberto ”, “ falta de informação sobre o controle da doença ” e “ presença de água parada ” ( sem ser específico quanto ao acondicionamento, tratamento ou limpeza ).
Foram exemplos de respostas incorretas : “ condições precárias nas edificações ” ; “ existência de terreno baldio ” ; “ falta de infraestrutura ” ; “ hábitos culturais inadequados ” ; “ alto grau de poluição ”, todos muito inespecíficos e sem possibilidade de pontuação.
Na resposta ao item ‘ b ’, uma boa parte dos alunos apresentou um dos exemplos do padrão de resposta, de medidas de prevenção e controle dos vetores a serem implementadas em residências.
Foram muito comuns as respostas “ usar areia em pratos de vasos de plantas ”, “ descartar e acondicionar pneus adequadamente ”, “ limpar periodicamente as calhas ” e “ não deixar água parada ” ( sem citar “ sem acondicionamento ou tratamento adequado ” ).
Alguns alunos citaram medidas de controle individual, como uso de repelentes corporais, e outros citaram ações governamentais que não são de uso residencial.
Foram exemplos de respostas incorretas : “ filtrar água ” ou “ tratar água ” ou “ tratar esgoto ” ; “ boas condições de higiene pessoal ” ( ou lavar as mãos ) ; “ limpar a residência ” ( muito inespecífico, sem possibilidade de pontuação ) ; “ usar repelente corporal ” ( é condição individual, não é condição a ser implementada em residências ) ; “ usar o fumacê ” ou “ usar fiscais sanitários ” ou “ educar ambientalmente ” ( são ações governamentais, não são condições a serem implantadas em residências ) ; e “ alto grau de poluição ” ( muito inespecífico, sem possibilidade de pontuação ).
Na resposta ao item ‘ c ’, uma parcela dos alunos apresentou um dos exemplos do padrão de resposta de ações governamentais que contribuem para a erradicação das doenças.
Foram muito comuns as respostas “ educação ambiental ”, “ promover ações de limpeza urbana ”, “ visita de agentes sanitários ” e “ multar os moradores que deixarem água parada com larvas ”.
Alguns alunos colocam deficiências governamentais como : “ falta de saneamento básico ” ou “ falta de informação ”.
Essas não são ações, mas, sim, deficiências e não podem ser pontuadas.
Foram exemplos de respostas incorretas : “ falta de uso de vacina ” ( a vacina ainda não foi desenvolvida e liberada ) ; “ tratamento de água ” ( muito inespecífico, no sistema de abastecimento de água as medidas de controle adequadas seriam controle de vazamentos e vedação dos reservatórios de distribuição ) ; e “ tratamento de esgoto ”.
As respostas dos estudantes revelaram um adequado conhecimento dos conceitos de epidemiologia e controle exigidos pela questão.
Uma vez que os comandos da questão solicitaram respostas diretas e objetivas, não havia necessidade de uma organização analítica, de capacidade de síntese ou competências ligadas à argumentação lógica.
Apesar disso, foi possível perceber -se grande dificuldade de redação do texto das respostas.
Naquelas avaliadas como fracas e medianas, os erros de conhecimento da Língua Portuguesa ( norma culta ) apontam para deficiências na estruturação e apresentação do pensamento por escrito. 3.3.2.3 - Análise da Questão Discursiva 4 do Componente de Conhecimento Específico A Tabela 3.19 contém as informações relativas à questão 4 do conjunto de questões do Componente de Conhecimento Específico.
O desempenho dos estudantes nessa questão foi inferior ao das questões 3 e superior ao da questão 5.
A Média geral do Brasil foi 44,3, sendo a maior Média registrada na região Nordeste ( 50,5 ), e a menor, na região Norte ( 38,0 ).
A nota Máxima ( 100,0 ) foi atingida por, pelo menos, um concluinte de todas as regiões.
A Mediana em todo o Brasil foi 40,0, o mesmo valor foi obtido nas regiões Norte, Sul e Centro-Oeste.
Nas demais regiões as medianas foram : 60,0 ( Nordeste ) e 50,0 ( Sudeste ).
As notas mínimas foram zero em todas as regiões.
O Gráfico 3.12 representa a distribuição de notas da questão discursiva 4, do Componente de Conhecimento Específico.
Essa distribuição é multimodal com moda principal nos alunos que deixaram a resposta da questão em branco, aproximadamente 17 \%.
Observam -se outros máximos locais no intervalo ( 50 ; 60 ], seguido de perto pelos intervalos ( 90 ; 100 ] e ( 30 ; 40 ], cada um com frequência em torno de 12 \%. 3.3.2.4 - Comentários sobre as respostas à Questão Discursiva 4 O enunciado da questão era claro e redigido em linguagem acessível.
Aborda o Estatuto da Cidade e os conceitos de Plano Diretor e de uso e ocupação do solo apresentados no conteúdo de Planejamento Urbano e Ambiental.
Aborda, ainda, conceitos de Impactos Socioambientais.
Dessa forma, a questão está adequada ao conteúdo previsto nas Diretrizes para Tecnologia em Gestão Ambiental da Portaria Inep nº 302, de 08 de junho de 2016.
Dentro dessa visão sistêmica, permite a avaliação de competências para planejar, executar e analisar instrumentos de gestão ambiental, políticas, programas e projetos, bem como aplicar a legislação pertinente à área.
O comando da questão solicitava que o estudante desenvolvesse uma resposta composta de duas partes, que objetivavam que o aluno apresentasse os seguintes conteúdos : três problemas socioambientais originados pelo adensamento/verticalização da cidade ; e um mecanismo que possibilite a participação da sociedade na gestão democrática, para a construção de uma cidade sustentável, explicando de que forma esse mecanismo dá voz e espaço aos diversos atores sociais.
Foram fornecidos um texto e duas figuras no enunciado que permitiam inferir, por interpretação, uma parte da resposta.
Em relação ao item ‘ a ’, os conhecimentos envolvidos apresentam baixo grau de subjetividade, o que levou as respostas a serem bem próximas do padrão de resposta.
Já no item ‘ b ’, que solicita mecanismo para dar voz e espaço aos atores sociais envolvidos, há um grau mais elevado de subjetividade, o que levou a uma maior diversidade de respostas.
Os comandos dos itens ‘ a ’ e ‘ b ’ eram claros, o que gerou poucas possibilidades de interpretações diferentes da esperada.
A linguagem estava adequada ao jargão inerente ao profissional da área de Tecnologia em Gestão Ambiental.
Em comparação com as demais questões de formação específica da Área, a questão 4 pode ser considerada de média dificuldade.
Nas respostas ao item ‘ a ’, houve baixa diversidade.
Para responder ao item ‘ b ’, alguns alunos parafrasearam trechos do enunciado, em tentativas de soluções quase sempre distantes do previsto no padrão.
O padrão de resposta contemplou uma variedade bem grande de possibilidades de resposta, fato que valorizou as diferentes abordagens escolhidas pelos estudantes ao formulá -las.
Assim, o padrão foi capaz de capturar as variações corretas de resposta.
As respostas que fugiram ao padrão estavam de fato incorretas.
Quanto ao uso da língua portuguesa, de forma geral, os alunos apresentaram muitos erros e, nem sempre, utilizaram termos específicos da Área.
Como esperado, observou -se uma maior aderência ao jargão específico da área entre os alunos com melhor desempenho na questão.
Nas respostas consideradas fracas, em geral, o estudante apenas citou um ou dois dos vários problemas socioambientais presentes no padrão de resposta ou ambos.
Os alunos que se enquadraram nessa faixa de desempenho, em geral, não formularam resposta para o comando que solicitava mecanismo de participação da sociedade no processo de construção da cidade sustentável.
Nas eventuais tentativas de resposta a esta solicitação, apresentaram deficiências de argumentação e textos desconexos, sem possibilidade de pontuação.
Convém destacar que o enunciado da questão cita as seguintes palavras do art. 2º do Estatuto da Cidade : (... ) “ gestão democrática por meio da participação da população e de associações representativas dos vários segmentos da comunidade na formulação, execução e acompanhamento de planos, programas e projetos de desenvolvimento urbano (... )".
Esse fragmento permitiu aos alunos mais articulados atender a parte da resposta esperada pelo padrão.
Nas respostas consideradas medianas, geralmente, o estudante citou dois ou três dos vários problemas socioambientais do padrão de resposta.
Alguns alunos, complementarmente, citaram um dos mecanismos de participação da sociedade válidos.
Com uma frequência menor, foram apresentadas respostas corretas para o comando que solicitava a forma do mecanismo de participação da sociedade no processo de construção da cidade sustentável.
Nas respostas consideradas boas, o estudante apresentou os três problemas socioambientais, citou o mecanismo por meio de um texto coerente e organizado para a explicação da forma de participação da sociedade no processo de construção da cidade sustentável.
Nessa faixa de desempenho, os candidatos demonstraram conhecer o jargão conceitual da área e possuir boa capacidade de argumentação.
Os eventuais deslizes nesses quesitos não comprometeram a solução apresentada.
No item ‘ a ’, deveriam ser apresentados três problemas socioambientais originados pelo adensamento e verticalização da cidade.
A maioria dos alunos citou algum dos itens do padrão.
Foram muito comuns as respostas “ necessidade de coleta, tratamento e destinação final de maior quantidade de resíduos ”, “ aumento da poluição das águas ”, e “ aumento da poluição ” ( sem especificar se da água, do solo, do ar, sonora ).
Alguns alunos citaram causas ou aspectos ambientais, ao invés de descrever consequências ou impactos : “ aumento populacional ” ou “ aumento na geração de resíduos ”.
Outros citaram causas ou aspectos ambientais, ao invés de descrever consequências ou impactos : “ aumento populacional ” ou “ aumento na geração de resíduos ”.
Alguns alunos citaram características que não têm alta relação de consequência com o adensamento e verticalização : “ aumento da criminalidade ” ou “ maior consumismo ” ou “ falta de gestão ” ou “ desigualdade social ”.
No item ‘ b ’, deveria ser citado um mecanismo que possibilitasse a participação da sociedade na gestão democrática para a construção de uma cidade sustentável, assim como a explicação da forma como que esse mecanismo dá voz e espaço aos diversos atores sociais.
Uma boa parte dos alunos apresentou um dos mecanismos do padrão de resposta, mas, nem sempre, soube explicar como esse mecanismo dá voz e espaço aos diversos atores sociais.
A maioria falou em “ audiências ou consultas públicas ”,"assembleias ou reuniões de bairro","através de ONGs", e poucos citaram “ orçamento participativo ”, “ conselho municipal ” ou"por plesbicito ou voto".
A explicação da forma com que esse mecanismo dá voz e espaço aos diversos atores sociais envolvidos quase sempre esteve ausente das respostas.
Uma parte citou “ opinando ” ou “ sugerindo escolha de solução ou alternativa ” ou “ se posicionando positiva ou negativamente sobre uma alternativa ”.
Outros, de forma incorreta, apontaram “ decidindo ” ( a população não decide, esse papel é dos Poderes Executivo ou Legislativo Municipal ).
De forma incorreta, alguns estudantes, ainda, citaram : “ no EIA/RIMA ” [ a população não tem voz no EIA/RIMA, esse é elaborado por grupo multidisciplinar contratado pelo empreendedor ( teria voz, apenas no caso da aplicação nesse processo do instrumento audiência ) ] ; “ por uma auditoria ambiental ” ( auditorias são realizadas pelo gestor para medir a efetividade de suas ações ambientais, a população não participa do processo ) ; ou “ através da educação ambiental ” ( ainda que desenvolva senso crítico nos estudantes, não é o mecanismo para participar das escolhas na construção das cidades ).
Outros poucos, de forma errada, ainda citaram medidas ou ações de engenharia visando ao controle e à sustentabilidade : “ reflorestamento ” ou “ controle do trânsito ” ( não foram pedidas medidas de intervenção ).
As respostas dos estudantes revelaram um adequado conhecimento dos conceitos de Plano Diretor e de uso e ocupação do solo apresentados no conteúdo de Planejamento Urbano e Ambiental.
Mostraram, ainda, conhecimento dos conceitos de Impactos Socioambientais apresentados.
De modo geral, as respostas demostraram que os estudantes têm muita dificuldade na expressão escrita do pensamento.
Os erros de coerência e coesão da argumentação além dos ortográficos, muitas vezes comprometem a compreensão do que o aluno buscou responder. 3.3.2.5 - Análise da Questão Discursiva 5 do Componente de Conhecimento Específico O desempenho dos estudantes de todo o Brasil nesta questão foi bem mais baixo do que nas outras duas questões discursivas desse componente.
A nota Média dos estudantes de todo o Brasil foi 14,0.
A maior Média foi registrada na região Nordeste ( 16,3 ), enquanto a menor Média foi registrada na região Centro-Oeste ( 11,4 ).
Quanto à variabilidade das notas, o Desvio padrão dos alunos do Brasil, como um todo, foi 24,3.
Enquanto o maior desvio foi encontrado na região Nordeste ( 25,2 ), o menor foi encontrado na região Centro-Oeste ( 23,3 ).
A Mediana para o Brasil e para todas as Grandes Regiões foi zero, indicando que mais do que 50 \% dos concluintes tiraram nota zero nessa questão e que a nota Mínima ( 0,0 ) foi a mesma em todas as regiões.
Para o conjunto de alunos de Tecnologia em Gestão Ambiental do Brasil, a nota Máxima foi 100,0, assim como em todas as regiões.
O Gráfico 3.13 representa a distribuição das notas da questão discursiva 5 do Componente de Conhecimento Específico.
Essa distribuição tem moda na nota zero.
Observa -se que a maior frequência de notas diferentes de zero ocorre no intervalo ( 10 ; 20 ]. 3.3.2.6 - Comentários sobre as respostas à Questão Discursiva 5 No que diz respeito ao item ‘ a ’, o comando estava claro e foi redigido em linguagem acessível.
Em relação ao item ‘ b ’, as respostas dos alunos evidenciaram que houve problema na interpretação do que foi solicitado.
A questão abordava a cadeia de operações do processo de reciclagem como um todo e o previsto na Política Nacional de Resíduos Sólidos, tema associado à Saúde e Saneamento Ambiental, mais especificamente, à gestão de resíduos sólidos : acondicionamento, coleta, transporte, tratamento e disposição final.
Dessa forma, a questão está adequada ao previsto nas Diretrizes para Tecnologia em Gestão Ambiental da Portaria Inep nº 302, de 08 de junho de 2016.
Dentro desse escopo, permitia avaliar a formação de competências para planejar, executar e analisar instrumentos de gestão ambiental, bem como aplicar a legislação pertinente à área.
O comando da questão solicitava que o estudante desenvolvesse uma resposta composta de duas partes : a primeira objetivava a escolha de dois elos da cadeia produtiva da reciclagem e a explicação para cada um deles e se a experiência relatada no estudo de caso contemplava o preconizado para o elo escolhido.
A segunda solicitava que o estudante informasse, justificando sua resposta, se a gestão adotada pela Prefeitura do estudo de caso, atendia o preconizado para as cooperativas de catadores constituídas por pessoas físicas de baixa renda.
No que diz respeito ao item ’ a ’, por conta de um razoável grau de subjetividade, o padrão era bastante amplo.
Trazia uma ou mais possibilidades de respostas para cada um dos elos que poderiam ser escolhidos pelos estudantes.
Em relação ao item ‘ b ’, foi comum os estudantes interpretarem a solicitação de citarem “ o papel atribuído às cooperativas citadas ” como sendo para citar “ o papel das cooperativas citadas ”.
Ou melhor, ficou uma dúvida sobre quem atribuiu o papel : a prefeitura ou a PNRS ( Política Nacional de Resíduos Sólidos ).
Assim, boa parte dos alunos respondeu apontando deficiências operacionais ou estruturais das cooperativas, que, na sua maioria, não são reguladas na PNRS.
Em comparação com as demais questões discursivas de formação específica, essa questão pode ser considerada difícil, pela necessidade de interpretação e análise do texto para responder aos comandos da questão.
O descompasso entre o que foi solicitado no item ‘ b ’ e a maneira como foi respondido, incorretamente pela maioria dos alunos, acabou por tornar o desempenho ainda pior do que o esperado nessa questão.
Principalmente no que tange ao item ‘ a ’, o padrão de resposta apresentou uma grande variedade de possibilidades, o que ajudou a pontuar quase todas as opções escolhidas pelos estudantes ao formular suas respostas.
Já no item ‘ b ’, houve um descompasso entre a resposta dada pela maioria dos alunos e o esperado no padrão de resposta.
No item ‘ a ’, a maioria dos alunos citou algum dos itens presentes no padrão.
Foi muito comum a escolha dos elos 1 e 3 com as respostas do tipo :"Elo 1 ( segregação na fonte ) – contemplado apenas parcialmente ( ou não totalmente ), pois a segregação é mal feita uma vez que as cooperativas reclamam que recebem muito resíduo orgânico misturado na coleta seletiva";"Elo 3 ( centrais de triagem ) – contemplado apenas parcialmente ( ou não totalmente ), pois as instalações físicas das cooperativas demandam adequação para recebimento e triagem"ou"demandam adequação para o armazenamento"ou"os estoques são acumulados com perda de espaço no pátio interno"ou"quatro das seis cooperativas não tem galpão próprio “.
Alguns citaram um elo e deram explicação com fragmento de texto relativo a outro elo.
Também houve os que confundiram a operação ou elo de reciclagem com o ciclo da reciclagem como um todo, respondendo, por exemplo, “ Elo 5 ( reciclagem ) não foi contemplado porque o percentual de recicláveis processado em relação ao total é baixo ”.
No item ‘ b ’, foi cobrado que o estudante explicasse se o papel atribuído às cooperativas atende às exigências da Lei nº 12.305/2010.
Como já comentado, houve um descompasso entre a resposta dada pela maioria dos alunos e o esperado no padrão de resposta, além do fato de que grande parte dos alunos não respondeu a esse item.
Foram muito comuns as respostas incorretas do tipo : “ atende apenas parcialmente, pois as instalações físicas das cooperativas demandam adequação para recebimento e triagem ” ; “ atende apenas parcialmente, pois as instalações físicas das cooperativas demandam adequação para o armazenamento ” ; “ atende apenas parcialmente, pois quatro das seis cooperativas não tem galpão próprio ” ; “ atende apenas parcialmente, pois comercialização é feita de forma independente pelas cooperativas, o que diminui o preço dos reciclados ” ; ou “ atende apenas parcialmente, pois as cooperativas tem a prática de descarte de recicláveis junto ao lixo convencional ”.
Outros poucos citaram :"não atende na medida em que as cooperativas praticam a queima a céu aberto de recicláveis de menor valor agregado, como papéis"; ou"não atende na medida em que a PNRS define uma ordem de prioridade na destinação dos resíduos ( não geração, redução na fonte, reutilização, reciclagem, tratamento e disposição ), que não foi respeitada, na medida em que, apenas uma parcela dos recicláveis foi processada".
Em um número bem menor foi citado que “ faltam incentivos ou investimentos financeiros sem dizer que a lei prioriza essa atuação por parte da prefeitura ”.
Um número muito baixo relatou que atende ( atende parcialmente ou não atende totalmente ) a PNRS, na medida em que essa incentiva a inclusão produtiva dos catadores através de cooperativas ( ou que essa prioriza a participação de cooperativas de pessoas físicas de baixa renda ou que prioriza o uso de cooperativas para tirar os catadores dos lixões.
Nas respostas consideradas fracas, em geral, o estudante confundiu conceitos ou escreveu resposta conceitualmente coerente, mas diferente da solicitada, por exemplo, confundiu a cadeia da reciclagem como um todo com cada uma das operações ( elos ) do processo.
Houve também respostas nas quais o aluno citou um elo, mas formula explicação relativa a outro.
Houve respostas genéricas, como as que citam elo 1 e elo 3, sem dar explicação de como a operação está relatada ( ou não ) no estudo de caso apresentado.
Em relação às respostas ao item ‘ b ’, a maioria dos alunos nem formulou resposta.
Existiram, ainda, aqueles que não responderam corretamente ao item ‘ b ’ e não ao item ‘ a ’.
Nas respostas consideradas medianas, o estudante respondeu, apenas, a um dos itens da questão deixando, na maioria das vezes, um segundo sem ser respondido ou com resposta incorreta.
Nas respostas consideradas boas, o estudante articulou corretamente as respostas, utilizando adequadamente o jargão conceitual da área e desenvolveu um texto estruturado, onde deslizes ocasionais não comprometeram o conteúdo apresentado.
Houve certa diversidade nas tentativas de resposta formuladas para o comando do item ‘ a ’.
Esse fato era esperado, dadas as diversas opções de escolha de elo e as várias características elencadas no estudo de caso.
Ainda assim, o padrão foi capaz de capturar as variações corretas de resposta.
No que diz respeito ao item ‘ b ’, o \textbf{elevado} número de respostas em branco ou de respostas mostrando deficiências operacionais ou estruturais das cooperativas, não reguladas na PNRS, denota problema de interpretação do comando do item, já comentado.
As respostas dos estudantes revelaram um baixo conhecimento dos conceitos relacionados à cadeia de operações do processo de reciclagem como um todo e o previsto na Política Nacional de Resíduos Sólidos.
O \textbf{elevado} número de respostas em branco, nota 0,0 ( zero ) ou 2,0 ( dois ) ratificam essa afirmação.
Como eram esperados na questão uma capacidade de análise e interpretação do estudo de caso e o desenvolvimento de explicações argumentativas, ficou demonstrado que, em geral, os estudantes têm muita dificuldade na expressão escrita do pensamento.
Não parecem consolidados, para a grande maioria dos concluintes avaliados, os conceitos relacionados à cadeia de operações do processo de reciclagem e o que é previsto na Política Nacional de Resíduos Sólidos. 3.3.3 - Considerações Finais Comparando as três questões da formação específica de Tecnologia em Gestão Ambiental, pode -se considerar a questão 3 fácil, a questão 4 mediana e a questão 5 muito difícil.
Possivelmente, os concluintes tinham pouca experiência com aplicação de conhecimentos teóricos a estudos de caso, experiência importante na formação dos profissionais de nível superior.
Tal experiência ajuda a ganhar habilidade de buscar soluções técnicas para problemas do cotidiano profissional.

% matched lemmas: elevado
\end{textsample}
